\documentclass[sigconf,balance=false]{acmart}
\usepackage{popets}

% Copyright
\setcopyright{popets}
\copyrightyear{YYYY}

% Issue info
\acmYear{YYYY}
\acmVolume{YYYY}
\acmNumber{X}
\acmDOI{XXXXXXX.XXXXXXX}
\acmISBN{}
\acmConference{Proceedings on Privacy Enhancing Technologies}
\settopmatter{printacmref=false,printccs=false,printfolios=true}

%%% CUSTOM IMPORTS
\usepackage{algpseudocode, algorithm, algorithmicx}

\newcommand*\Let[2]{\State #1 $\gets$ #2}
\newcommand{\share}[1]{[\![#1]\!]}
\DeclareMathOperator{\nnz}{nnz}

\begin{document}

%%
%% The "title" command has an optional parameter,
%% allowing the author to define a "short title" to be used in page headers.
\title[Sparsified Secure Aggregation]{Sparsified Secure Aggregation\\ and its Applications to Large-Scale Federated Learning}

%%%%%%%%%%%%%%%% Authors' Info %%%%%%%%%%%%%%%%%
%%
%% The "author" command and its associated commands are used to define
%% the authors and their affiliations.

\author{Marc Damie}
\orcid{0000-0002-9484-4460}
\affiliation{%
	\institution{University of Twente}
	\city{Enschede}
	\country{The Netherlands}}
\email{m.f.d.damie@utwente.nl}


%%
%% By default, the full list of authors will be used in the page
%% headers. Often, this list is too long, and will overlap
%% other information printed in the page headers. This command allows
%% the author to define a more concise list
%% of authors' names for this purpose.

\renewcommand{\shortauthors}{Damie}

%%
%% The abstract is a short summary of the work to be presented in the
%% article.
\begin{abstract}
	A clear and well-documented
\end{abstract}


%%
%% Keywords. The author(s) should pick words that accurately describe
%% the work being presented. Separate the keywords with commas.
\keywords{Secure Aggregation, Secret Sharing, Sparsity, Differential Privacy}


\maketitle

\section{Introduction}

\section{Background}
\subsection{Federated Learning}

\subsection{Gradient Sparsification}

\subsection{Secret Sharing}

\subsection{Differential Privacy}

\subsection{Threat Model}

\section{Problem statement}

This section proposes to circumvent the problem of secure sparse aggregation and focus on a variant: probabilistic sparse agreggation.
This alternative would support the same ML use cases (especially sparse gradient aggregation), but has properties that could simplify the design of secure protocols.

\paragraph{Definition}

Let $y_*$ be the output of an exact sparse aggregation. A probabilistic sparse aggregator outputs $\widetilde{y}$ such that $\mathbb{E}(\widetilde{y}) = y_*$.
This probabilistic output seems acceptable for ML applications (e.g., decentralized gradient descent).

For simplicity, we assume a probabilistic sparse aggregator only outputs sparse vectors with one non-zero elements.
However, we can naturally extend this definition to any probabilistic ouput with $k$ non-zero elements.

The probabilistic sparse aggregation would be executed several times in a row to approach the perfect sparse aggregation results.

\paragraph{Computational advantages}
This probabilistic problem might have more efficient solutions than the classic sparse aggregation.
We may build protocols with logarithmic communication costs since the protocols must only output one non-zero value.
Hence, we would not have a cost linear in the number of non-zeros (or even linear in the number of dimensions).

\paragraph{Security advantages}
This problem also solves our leakage issue.
Indeed, as the protocol must output one non-zero elements, we do not need to leak securely the number of non-zeros (contrary to the exact sparse aggregation).

\section{Sparsified secure aggregation protocol}

\subsection{Threat Model}


\subsection{Base protocol}
This section presents a secure probabilistic sparse aggregation.
We rely on secret-sharing-based multi-party computation and assume that all the data owners are involved in the computation.
We have $N$ parties each with a $d$-dimensional sparse vector $x_k$ containing a single non-zero elements $(i_k, v_k)$.

\paragraph{Toy protocol}

A first simplistic solution would be to agree on a public random seed.
The parties would use this random seed to sample an index $i$ in $\{1\dots d\}$.
Then, each party shares a secret value $s_k$ with the rest of the parties: $s_k = v_k$ if $i = i_k$, $0$ otherwise.
The parties sum all the shares and reveal the sum multiplied by $d$.

This probabilistic aggregation is unbiased, but it has a slow convergence.
Indeed, picking a coordinate at random in sparse vectors will result in picking zero coordinates most of the time.
Hence, the protocol is efficient but would require many repetitions to have an accurate estimation of the perfect sparse aggregation.
To improve the previous protocol, we need to sample a coordinate \emph{only among the non-zero coordinates} instead of sampling it in $\{1\dots d\}$.
Such improvements would guarantee a faster convergence of our estimator.

\paragraph{Biased estimation}

Algorithm \ref{alg:prob_sparse_agg} provides a first efficient solution to our problem.
This iterative protocol samples one coordinate among the non-zero coordinates.
The communication cost (per data owner) is $O(N\cdot \log d)$ (for $N$ parties and $d$ dimensions).
Our sampling protocol represents the coordinate space (i.e., $\{1\dots d\}$) via a binary tree, and samples `a non-zero path' in the tree.

\begin{algorithm}
	\caption{Biased probabilistic sparse aggregation \label{alg:prob_sparse_agg}}

	\begin{algorithmic}[1]
		\Require{A group of $N$ parties. Each party $k$ owns a private sparse $d$-dimensional vector $x_k$.}
		\Statex
		\Function{ProbabilisticSparseAggregation}{$\{x_k\}_{k\in[N]}$}
		\Let{$l$}{$1$}
		\Let{$u$}{$d$}
		\While{$l \neq u$}
		\State{// Local operations:}
		\Let{$m$}{$\lfloor\frac{u+l}{2}\rfloor$}
		\Let{$c^\text{left}_k$}{Number of nnz in $x_k$ with a coordinate in $\{l\dots m\}$ in $x_k$}
		\Let{$c^\text{right}_k$}{Number of nnz in $x_k$ with a coordinate in $\{m+1\dots u\}$ in $x_k$}
		\Statex
		\State{// Online operations:}
		\For{all $i\in [N]$}
		\Let{$\share{c^\text{left}_k}$}{Party $k$ shares $c^\text{left}_k$ with the other parties}
		\Let{$\share{c^\text{right}_k}$}{Party $k$ shares $c^\text{right}_k$ with the other parties}
		\EndFor
		\Statex
		\Let{$\share{c^\text{left}}$}{$\sum_{k\in[N]}\share{c^\text{left}_k}$}
		\Let{$\share{c^\text{right}}$}{$\sum_{k\in[N]}\share{c^\text{right}_k}$}
		\State{Generate two random secret-shared numbers: $\share{r_1}, \share{r_2}$}

		\Statex
		\State{// Choose uniformly at random a path \textit{containing non-zero coord.}:}
		\Let{$\text{choice}$}{\Call{Reveal}{$\share{c^\text{left}}\cdot \share{r_1}< \share{r_2}\cdot\share{c^\text{right}}$}}
		\If{$\text{choice}$}
		\Let{$l$}{$m$}
		\Else
		\Let{$u$}{$m$}
		\EndIf
		\EndWhile
		\Statex
		\Let{$\text{coord}$}{$l$}
		\For{all $i\in [N]$}
		\Let{$\share{\text{val}_k}$}{Party $k$ shares the value of $x_k$ on coordinate $\text{coord}$}
		\EndFor
		\Let{$\text{val}$}{\Call{Reveal}{$\sum_{k\in[N]} \share{\text{val}_k}$}}
		\State \Return{$(\text{coord}, \text{val})$}
		\EndFunction
	\end{algorithmic}
\end{algorithm}

However, this protocol outputs a biased estimation of the exact sparse aggregation.
The causes of this bias are: (1) the tree path sampling is not uniform, (2) even with a uniform sampling, the expected value would be $\mathbb{E}(\widetilde{y}) = y_*\cdot \nnz_*$, with $nnz_*$ the number of non-zeros in the exact output $y_*$.


\paragraph{Unbiased estimation}

To have a uniform path sampling, we can add a simple permutation step.
At the beginning of the protocol, the parties agree on a public random seed and use it to generate a permutation $\pi$ applied to the non-zero coordinates.
Every time the protocol is executed, a new permutation is generated.
Thanks to this random permutation, the algorithm would perform a uniform sampling over the non-zero coordinates of the aggregated vector.

The second source of bias is less trivial to solve because we need to estimate $nnz_*$ without computing the exact sparse aggregation.
For this estimation, we will rely on the following proportion: $\rho = \frac{{\sum_k \nnz(x_k)}}{\nnz_*}$.
This proportion can be interpreted as an overlapping rate, i.e., the average number of values aggregated on non-zero coordinates.
In other words, this value quantifies how much the input sparse vectors overlap.

To estimate $\rho$, we can also use our probabilistic aggregation protocol to sample multiple non-zero coordinates and compute the average number of non-zeros on the sampled coordinates.
Instead of sharing values $\share{\text{val}_k}$, the parties would share either $1$ (if they have a non-zero on the sampled coordinate), $O$ otherwise.
This estimation protocol can reuse the tree paths already sampled for the probabilistic sparse aggregation.
Once we have an estimation of $\rho$, we can deduce the number of non-zeros in the complete aggregated vector.
Note that $\rho$ can remain secret because this estimation protocol can be performed under secret sharing.
Hence, we can output directly an unbiased probabilistic sparse aggregation (without leaking the biased result and the correction factor).

\subsection{Differentially private protocol}

\section{Theoretical Analysis}

\subsection{Differential Privacy}

\subsection{Convergence}

\subsection{Communication Cost}

\section{Evaluation}

\subsection{Utility}

\subsection{Efficiency}

\section{Extensions}
\subsection{Malicious Security}

\subsection{Single-server variant}

\section{Conclusion}

%%
%% The acknowledgments section is defined using the "acks" environment
%% (and NOT an unnumbered section). This ensures the proper
%% identification of the section in the article metadata, and the
%% consistent spelling of the heading.
% \begin{acks}
% 	To Robert, for the bagels and explaining CMYK and color spaces.
% \end{acks}

%%
%% The next two lines define the bibliography style to be used, and
%% the bibliography file.
\bibliographystyle{ACM-Reference-Format}
\bibliography{references}


%%
%% If your work has an appendix, this is the place to put it.
% \appendix

\end{document}
\endinput
%%
%% End of file `sample-sigconf.tex'.
